% ============================================================
%  黄江南 · 软件开发工程师（Go / AI Agent）· 2027 届
%  resume4：删滴滴 / 删创业 / 删顶部标签栏
%           项目 = AegisNest（客服）+ MockAtlas（面试）
%           徽标仅 Apache / 小红书；Casdoor、南邮无徽标
%
%  编译：.\build.ps1   （一次产出 彩色 + 黑白 两个 PDF）
%  彩色/黑白切换：把下面 \bwfalse 的 false 改成 true
%  徽标：..\logos\<name>-color.png / <name>-gray.png（共享目录）
% ============================================================
\newif\ifbw
\bwfalse % ================= TOGGLE: false = 彩色版，true = 黑白版 =================

\documentclass[10pt,a4paper]{awesome-cv-zh}

\usepackage{graphicx}
\graphicspath{{logos/}{../logos/}}

\geometry{left=1.10cm,right=1.10cm,top=0.64cm,bottom=0.45cm}

% ---------------- 配色 ----------------
\definecolor{accentblue}{HTML}{1F4E79}
\ifbw
  \colorlet{accent}{black}
  \colorlet{sectiondivider}{black}
  \colorlet{awesome}{black}
  \colorlet{text}{black}
  \colorlet{graytext}{gray}
  \colorlet{accenttint}{black!7}
\else
  \colorlet{accent}{accentblue}
  \colorlet{sectiondivider}{accentblue}
  \colorlet{awesome}{accentblue}
  \colorlet{accenttint}{accentblue!9}
\fi

\pagestyle{empty}
\setlength{\parindent}{0pt}
\setlength{\parskip}{0pt}

% ---------------- 徽标 ----------------
\ifbw
  \newcommand{\badge}[1]{\raisebox{-0.16em}{\includegraphics[height=1.12em]{#1-gray.png}}\hspace{1.0mm}}
\else
  \newcommand{\badge}[1]{\raisebox{-0.16em}{\includegraphics[height=1.12em]{#1-color.png}}\hspace{1.0mm}}
\fi

% ---------------- 区块标题 ----------------
\newcommand{\rsection}[1]{%
  \par\addvspace{1.15mm}%
  \noindent{\color{accent}\rule[0.10ex]{1.6mm}{3.0mm}}\hspace{1.6mm}{\fontsize{11pt}{12pt}\selectfont\bfseries\color{text}#1}\par
  \vspace{0.5mm}%
  \noindent{\color{sectiondivider}\rule{\textwidth}{0.6pt}}\par
  \vspace{0.5mm}%
}

% ---------------- 条目头 ----------------
\newcommand{\entry}[3]{%
  \par\addvspace{0.6mm}%
  \noindent{\fontsize{9.7pt}{11.5pt}\selectfont\bfseries\color{text}#1}\hspace{1.6mm}{\fontsize{8.8pt}{10.6pt}\selectfont\bfseries\color{accent}#2}\hfill{\fontsize{8.3pt}{10pt}\selectfont\color{graytext}#3}\par
  \vspace{0.3mm}%
}

% ---------------- 列表 ----------------
\newenvironment{items}{%
  \begin{itemize}[leftmargin=1.10em,label={\color{accent}\scriptsize\textbullet},topsep=0.6pt,itemsep=0.8pt,parsep=0pt,partopsep=0pt]%
}{\end{itemize}}

\begin{document}
\fontsize{8.9pt}{10.0pt}\selectfont

% ======================= 头部 =======================
\noindent
\begin{minipage}[b]{0.62\textwidth}
  {\fontsize{23pt}{25pt}\selectfont\bfseries\color{text} 黄江南}\\[1.5pt]
  {\fontsize{10.3pt}{12pt}\selectfont\bfseries\color{accent} 软件开发工程师 · 后端（Go）/ AI Agent}\\[1.5pt]
  {\fontsize{8.4pt}{10pt}\selectfont\color{graytext} 2027 届 · 南京邮电大学 · 求职意向：后端开发 / Agent 开发}
\end{minipage}%
\hfill
\begin{minipage}[b]{0.36\textwidth}\raggedleft
  {\fontsize{8.4pt}{11.6pt}\selectfont\color{text}%
  3310145612@qq.com\\GitHub: JasonBuildAI \;·\; 微信: HJN522312MOST}
\end{minipage}
\par\vspace{1.2pt}
\noindent{\color{accent}\rule{\textwidth}{1.2pt}}\par
\vspace{0.4pt}

% ======================= 教育背景 =======================
\rsection{教育背景}
\noindent{\fontsize{9.6pt}{11.2pt}\selectfont\bfseries\color{text} 南京邮电大学（双一流）}\hspace{1.6mm}{\fontsize{8.9pt}{10.8pt}\selectfont\bfseries\color{accent} 微电子科学与工程 · 本科}\hfill{\fontsize{8.3pt}{10pt}\selectfont\color{graytext} 2023.09 -- 2027.06}\par
{\fontsize{8.5pt}{10.4pt}\selectfont CET-6 / CET-4 \;·\; \textbf{核心基础}：计算机网络 / 高级语言程序设计 / 操作系统 / 数据库}\par

% ======================= 实习经历 =======================
\rsection{实习经历}
\entry{\badge{apache}Apache Casbin（Apache 在孵项目）· Casbin 创始团队}{2026 明日之星预选生}{2026.04 -- 2026.06}
{\fontsize{8.5pt}{10.4pt}\selectfont\color{graytext}%
导师：hsluoyz（北京大学罗杨教授，Apache Casbin 社区负责人）\;·\; 项目：\textbf{OpenAgent}（\texttt{the-open-agent/\allowbreak openagent}）—— Casbin 原班人马孵化的独立开源项目，GitHub 5,600+ Stars；本人为\textbf{官方 Member}，\textbf{累计代码变更超 1 万行}，撰写技术文章累计浏览超 5 万、带动项目 Star 增长 1,500+。}\par
\vspace{0.5mm}
{\fontsize{8.5pt}{10.4pt}\selectfont%
\textbf{项目背景}：OpenAgent 是面向个人本地部署的开源 AI Agent，采用 Go 单进程架构，将 Agent 运行时、工具调用、多模型接入与 Web 控制台收敛为一个预编译二进制：\textbf{单文件约 23 MB、冷启动约 2.7 s}，相较同类 Node.js 方案（29.2 s、4.2 万+ 文件、382 MB）\textbf{快约 11 倍、体积约为其 1/16}。}\par
\vspace{0.4mm}
\begin{items}
  \item \textbf{Windows GUI 自动化引擎}：\textbf{从 0 到 1 独立设计并交付}基于 UIA 的 GUI 工具链，新增 UIA 引擎、任务调度器与 9 个原子工具（元素查找 / 交互 / 等待 / 断言 / 应用启动 / 系统指标读取等），并同步提供非 Windows 平台降级实现与安全熔断。
  \item \textbf{AgentBench 评测框架}：\textbf{从 0 到 1 独立设计并交付} Agent 能力基准框架 \texttt{the-open-agent/agentbench}（Python 核心 1692 行 / 33 文件），覆盖 \textbf{8 大套件、104 个任务}，内置 bootstrap 置信区间统计，评测可脚本化复现；据此完成对标 OpenClaw 的横向实测：工具调用成功率 \textbf{83.33\% vs 61.11\%}、单次耗时 \textbf{14.2s vs 26.1s}、Token 消耗 \textbf{1312 vs 2362}。
  \item \textbf{会话标题容错链路}：设计并实现会话标题自动生成与降级链路，通过 \texttt{ResolveChatTitle} 等函数完成用户消息兜底、HTML 清洗、空白折叠与 rune 安全截断，保证模型调用失败时标题仍可用。
  \item \textbf{性能、部署与大型特性}：移除 \texttt{GetAnswer} 接口重复 LLM 调用、补齐 Casdoor 不可用时的默认用量元数据、对齐 docker-compose 端口与数据库主机；推进动态模型目录与自适应定价（覆盖 15 个 Provider）与站点知识自动生成 Skill 两项大型特性。
\end{items}

\entry{Casdoor（CNCF 云原生全景图收录 · Go + React · 14,000+ Stars）}{开源贡献者}{2026.04}
\begin{items}
  \item 重构并改进项目 README（已合入），新增简体中文 README 与语言切换入口，优化文档链接与章节结构；并参与社区运营与对外推广，提升中文开发者的上手体验与文档可维护性。
\end{items}

% ======================= 开源经历 =======================
\rsection{开源经历}
\entry{\badge{xiaohongshu}小红书 FireRed-OpenStoryline（Python + JavaScript · 3,385 Stars / 394 Forks）}{核心贡献者}{2026.02 -- 2026.04}
{\fontsize{8.5pt}{10.4pt}\selectfont\color{graytext}%
\textbf{项目背景}：小红书 FireRedTeam 开源的 AI 视频剪辑 Agent，以自然语言交互 + LLM 规划 + 工具编排把手工剪辑转变为“意图驱动”的创作流程；本人在所有贡献者中 \textbf{PR 数第二}，独立承担会话状态持久化与恢复、会话历史、媒体链路与可用性治理等主干模块。}\par
\vspace{0.4mm}
\begin{items}
  \item \textbf{会话状态持久化与安全恢复}：独立设计并实现会话状态落盘与恢复机制——对损坏的工具调用行采取 fail-closed 策略而非静默跳过；为会话历史与 LangChain 消息分别设置 2000 / 4000 条上限；构建递归密钥脱敏链路，按字段名与内联赋值两种模式识别 \texttt{client\_secret} / \texttt{refresh\_token} / \texttt{x-api-key} / \texttt{access\_key} 等敏感字段并统一替换为 \texttt{***}；对 \texttt{getSession} 异常做类型化处理。
  \item \textbf{会话历史与媒体链路}：独立完成 Web 端对话历史侧栏（面板 DOM / 样式 / 中英双语 i18n / \texttt{aria} 无障碍标签），基于 localStorage 实现会话 ID 与列表双键存储、历史会话渲染切换与切换保护；后端抽象 MCP 媒体传输策略 \texttt{inline\_media = auto / always / never}，按 \texttt{connect\_host} 自动判定同机与跨机链路，同机路径直传、跨机才内联 base64，避免远程模式下 base64 放大导致的超大 payload。
  \item \textbf{Agent 可用性与稳定性治理}：新增 API Key 前置校验，任务开始前向 OpenAI 兼容接口发一次最小请求（\texttt{max\_tokens=1} 控制成本、10s 超时），区分 401 / 404 / 429 / 5xx / 超时 / 连接失败并返回可操作提示；为媒体搜索与下载补齐指数退避重试（初始 1s、2 倍退避、上限 10s），区分网络错误与非网络错误；统一日志出口并完善媒体文件保存的错误处理。
\end{items}

% ======================= 项目经历 =======================
\rsection{项目经历}
\entry{AegisNest 智能客服系统（Python + FastAPI + LangGraph）}{个人项目}{2026.07 -- 2026.09}
{\fontsize{8.5pt}{10.4pt}\selectfont\color{graytext}%
\textbf{项目介绍}：面向 API 开放平台的 Multi-Agent 智能客服系统，以 LangGraph 的 DAG 编排协作为核心，集成意图路由、RAG 混合检索、工具诊断、会话记忆、安全脱敏、多级兜底等能力，为开发者提供全流程智能技术支持。}\par
\vspace{0.4mm}
\begin{items}
  \item \textbf{Multi-Agent 编排}：基于 LangGraph StateGraph 设计意图路由、文档 RAG、API 诊断、账单、工单、安全审查 6 个专业 Agent 的职责划分与 DAG 管线，由中心编排器统一调度，安全审查节点串行收口，保证安全审查一定在最终输出之前执行。
  \item \textbf{RAG 混合检索}：实现向量检索（text-embedding-v3，1024 维）+ BM25 双路并行召回，经 RRF 融合、gte-rerank-v2 精排、上下文压缩后带引用生成；为错误码类问题设计标量热路径直查数据库，该环节耗时\textbf{从 0.74s 降到 0.010s（约 78 倍）}。
  \item \textbf{性能优化}：设计语义缓存、路由缓存与模型分层，叠加错误码热路径和 Multi-Agent 并行策略，冷启动 P95 \textbf{从 48.58s 降到 6.83s（约 7 倍）}，吞吐\textbf{从 0.18 提升到 0.75 req/s}，热点问题命中缓存仅 \textbf{0.38s、零 token 消耗（约 26 倍加速）}。
  \item \textbf{安全、兜底与可观测}：实现覆盖 API Key、Token、手机号等八类敏感信息的三层脱敏（用户输入、工具日志、最终输出），安全审查节点在管线末端收口；多级兜底覆盖从单 Agent 异常到整条管线崩溃的各种失败；基于逐节点链路追踪实现全链路可观测与按租户成本统计。
\end{items}

\entry{MockAtlas 智能面试助手（Go + Eino · Multi-Agent）}{个人项目}{2025.12 -- 2026.02}
{\fontsize{8.5pt}{10.4pt}\selectfont\color{graytext}%
\textbf{项目介绍}：基于 Go 与 Eino 框架的 AI 智能面试助手，以 Multi-Agent 的 DAG 编排协作为核心，集成 RAG 多路召回题库检索、动态难度调节、Agent 记忆系统等能力，提供简历分析、智能出题、实时面试、评估报告的全流程模拟面试体验。}\par
\vspace{0.4mm}
\begin{items}
  \item \textbf{Multi-Agent 编排}：设计 7 个 Agent 的职责划分与 7 阶段 DAG 编排流程，采用中央编排器模式管理 Agent 间数据流转，实现从 JD 分析、简历匹配、出题规划到面试评估的全链路自动化。
  \item \textbf{RAG 多路召回}：实现 Milvus 向量检索（1024 维）+ 纯内存 BM25 关键词检索双路并行召回，通过 RRF(k=60) 融合排序 + LLM Rerank 重排序输出最终候选题，支持按 \texttt{user\_id} 字段的用户级数据隔离。
  \item \textbf{记忆、难度调节与 Skill}：实现滑动窗口短期记忆 + Redis/MySQL 长期用户画像存储与薄弱知识点自动追踪；采用两阶段出题 + 动态难度调节 + 智能追问实现难度自适应；设计有状态多轮交互 Skill 接口（区别于无状态 Tool），开发快速测验、概念教学等 4 个内置 Skill，通过优先级匹配的 SkillRegistry 统一管理。
  \item \textbf{RAG 质量评估与调优}：构建离线评估流水线，基于 50 条标注样本计算 Recall@K / MRR 指标并按 topic 分组统计；通过 A/B 对比将 TopK 从 10 调至 20，\textbf{Recall@10 提升 7.3\%}、\textbf{MySQL 领域召回率提升 16.7\%}。
\end{items}

% ======================= 个人技术栈 =======================
\rsection{个人技术栈}
\begin{items}
  \item \textbf{技术栈}：Go、Python、TypeScript / JavaScript、C；MySQL、Redis、消息队列、RESTful API；LangChain / LangGraph / Eino、RAG（向量检索 / BM25 / RRF / Rerank）、Tool Calling、MCP / A2A；Docker、Linux、Git、React / Vue。
\end{items}

\end{document}
